% Recuerda incluir en el preámbulo:
% \usetikzlibrary{arrows.meta, babel}

\begin{tikzpicture}[
  scale=0.85, 
  auto=left,
  >={Stealth[length=2.5mm]},
  every node/.style={circle, fill=blue!20, draw=black, line width=1pt, inner sep=2pt, minimum size=6mm},
  flecha/.style={->, line width=1pt}
]

  % Vértices según la disposición del dibujo
  \node (v1) at (2,4) {$v_1$};
  \node (v2) at (0,2) {$v_2$};
  \node (v3) at (4,2) {$v_3$};
  \node (v4) at (2,0) {$v_4$};
  \node (v5) at (6,4) {$v_5$};
  \node (v6) at (8,2) {$v_6$};
  \node (v7) at (6,0) {$v_7$};

  % Aristas orientadas del bloque izquierdo
  \draw[flecha] (v2) -- (v1);
  \draw[flecha] (v1) -- (v3);
  \draw[flecha] (v2) -- (v3);
  \draw[flecha] (v2) -- (v4);
  \draw[flecha] (v4) -- (v3);

  % Conexión puente orientada
  \draw[flecha] (v3) -- (v5);

  % Aristas orientadas (v6 -> v5 invertida)
  \draw[flecha] (v6) -- (v5);
  \draw[flecha] (v5) -- (v7);

  \node[draw=none, fill=none, scale=1.5] at (4,-1.5) {$D$};
\end{tikzpicture}
